\section{Additive recurrence and the exclusion of simple colors}
\label{sec:recurrence}

We retain the odd division degree hypothesis and the hypothetical quotient
\(\phi:U(k)\to\mathcal F\) from Section~\ref{sec:palettes}.  Our objective is
to show that membership in \(\mathcal S\) is unchanged by scalar translations
on the additive chart.  Once this is known, fractional transformations
generate left rotations whose colors generate \(\mathcal F\), contradicting
the fact that \(\mathcal S\) is nonempty and proper.  This follows the
organization of \cite[Sections~5.5--5.6]{MPNF}.  The recurrence proofs below
are specific to positive characteristic: odd characteristic permits a
quadratic Fourier argument, whereas characteristic two requires additive
polynomials and the subgroup-complement property of
Theorem~\ref{fin:obstruction}.

Write \(\tau_s\) for the scalar translation \(x\mapsto x+s\), with
\(s\in k\).  For a palette \(c\in\mathcal C\), put
\[
 I_c(h)=\boldsymbol 1_{\mathcal S}(c_h),\qquad h\in\mathcal H.
\]
We shall produce a palette satisfying
\begin{equation}\label{rec:translation-invariance}
 I_c(\tau_s h)=I_c(h)\qquad(s\in k,\ h\in\mathcal H).
\end{equation}
Recall that the nonnegative function \(K\) constructed in
Section~\ref{sec:palettes} has a spectral representation
\[
 K(b)=\int_{\widehat J}\chi(b)\,d\sigma(\chi),\qquad
 \sigma(\widehat J)=\delta,\qquad \sigma(\{1\})\geq\delta^2,
 \qquad \delta=|\mathcal F|^{-|\mathcal F|}.
\]
Here additive groups have the discrete topology, and their duals are
compact groups of characters with values in the complex unit circle.

\subsection{Two averaging facts}

We first isolate the Fourier argument shared by the two characteristics.
An average on a countable set \(\Omega\) will mean integration against a
finitely supported probability measure on \(\Omega\), denoted by
\(\mathbb E_j\).  A subset has density zero for these averages if its
indicator has average tending to zero.

\begin{lemma}\label{rec:spectral-concentration}
Let \(G\) be a countable abelian group, let \(v:\Omega\to G\), and suppose
that
\begin{equation}\label{rec:orthogonality}
 \mathbb E_j\chi(v(y))\longrightarrow
 \begin{cases}1,&\chi=1,\\0,&\chi\ne1\end{cases}
 \qquad(\chi\in\widehat G).
\end{equation}
Let \(E\subseteq\Omega\).  For each \(x\notin E\), suppose that there are
\(a_x\in G\) and a positive measure \(\nu_x\) on \(\widehat G\), of mass
\(\delta\) and with \(\nu_x(\{1\})\geq\delta^2\), such that
\[
 F_x(y)=\int\chi(v(y)-a_x)\,d\nu_x(\chi)=0
 \qquad(y\in E).
\]
If \(\Omega\setminus E\) is nonempty, its lower density is at least
\(\delta\).  If, along a subsequence, \(E\) has density tending to
\(d>0\), then along a further subsequence there is a finite set
\(\Lambda\subset\widehat G\setminus\{1\}\) such that
\begin{equation}\label{rec:finite-spectral-list}
 \nu_x(\Lambda)\geq d\delta^2/2\qquad(x\notin E).
\end{equation}
\end{lemma}

\begin{proof}
By dominated convergence and \eqref{rec:orthogonality},
\(\mathbb E_jF_x\to\nu_x(\{1\})\).  Since \(|F_x|\leq\delta\) and
\(F_x\) vanishes on \(E\), this gives the lower-density assertion.

The union of the atoms of the measures \(\nu_x\) is countable.  Pass to
a subsequence on which
\[
 c(\chi)=\lim_j\mathbb E_j
       \bigl(\boldsymbol 1_E(y)\chi(v(y))\bigr)
\]
exists for every character in that union.  In particular, \(c(1)=d\).
For any finite list of distinct characters, their Gram matrices for
\(\mathbb E_j\) converge to the identity by
\eqref{rec:orthogonality}.  Bessel's inequality therefore gives
\(\sum_\chi|c(\chi)|^2\leq1\).  Consequently the set of nontrivial atoms
for which \(|c(\chi)|\geq d\delta/2\) is finite; denote it by
\(\Lambda\).

The nonatomic part of \(\nu_x\) contributes zero to the limiting average
against \(\boldsymbol 1_E\).  Indeed, the mean square of its Fourier
transform, evaluated at \(v(y)-a_x\), tends to zero: expand the square,
apply \eqref{rec:orthogonality} to \(\chi\overline\eta\), and integrate
over two copies of the nonatomic measure.  The diagonal has product
measure zero.  Cauchy--Schwarz now proves the assertion about the weighted
average.  The atomic part may be averaged term by term because its
weights are summable.  Thus \(F_x|_E=0\) implies
\[
 0=d\nu_x(\{1\})+
   \sum_{\chi\ne1}c(\chi)\chi(-a_x)\nu_x(\{\chi\}).
\]
The absolute contribution from characters outside \(\Lambda\cup\{1\}\)
is at most \(d\delta^2/2\).  The trivial character contributes at least
\(d\delta^2\), and \(|c(\chi)|\leq1\).  The remaining atoms must
therefore have total mass at least \(d\delta^2/2\), as required.
\end{proof}

We shall also use the following elementary form of polynomial avoidance.
It applies to finite additive subgroups without any compatibility with
the coordinates of the ambient vector space.

\begin{lemma}\label{rec:polynomial-density}
Let \(k\) be an infinite field of positive characteristic.  A proper
algebraic subset of \(k^r\) has density zero on every increasing sequence
of finite additive subgroups exhausting \(k^r\).  It also has density
zero on products of finite subsets whose smallest cardinality tends to
infinity.
\end{lemma}

\begin{proof}
It suffices to consider the zero set of a nonzero polynomial \(f\) of
total degree \(D\).  The elementary grid bound, obtained by induction
from the one-variable root bound, says that its zero proportion on
\(A_1\times\cdots\times A_r\) is at most
\(D/\min_i|A_i|\).  This proves the second assertion.

For the first assertion, choose any finite additive subgroup
\(H\subset k\).  An exhausting sequence \(U_j\subset k^r\) eventually
contains \(H^r\), and each such \(U_j\) is a disjoint union of its
\(H^r\)-cosets.  On a coset \(u+H^r\), apply the grid bound to the
nonzero polynomial \(f(u+X)\).  The zero proportion is at most
\(D/|H|\), uniformly in the coset.  Finite additive subgroups of \(k\)
have arbitrarily large cardinality, so this bound proves the claim.
\end{proof}

\subsection{Odd characteristic}

Assume first that \(\operatorname{char}k\ne2\).  Our choice of Cayley
parameter gives \(\bar\theta=-\theta\), so that \(T=0\).  Let
\(\mathcal A=k^\times\ltimes k\) be the scalar affine group, with
\((l,s)\) acting by \(z\mapsto lz+s\).  Associate to its elements the
vectors
\begin{equation}\label{rec:quadric}
 v(l,s)=\left(l,s,\frac{s^2+P_\theta}{l}\right)\in k^3.
\end{equation}
For a fixed initial affine position, the scalar parameters in
Equation~\ref{pal:scalar-parameters} become linear functions of these
three coordinates.  We average over \(l\in I_j\setminus\{0\}\) and \(s\in H_j\)
independently and uniformly, where \(I_j,H_j\) are increasing finite
additive subgroups exhausting \(k\), and
\[
 \{ab:a,b\in I_j\}\subseteq H_j.
\]
Such sequences exist by taking finite-dimensional vector spaces over
the prime field and enlarging \(H_j\) as necessary.

\begin{lemma}\label{rec:quadric-orthogonality}
For every nontrivial additive character \(\chi\) of \(k^3\),
\(\mathbb E_j\chi(v(l,s))\to0\).  Every proper polynomial locus in
\((l,s)\) also has density zero for these averages.
\end{lemma}

\begin{proof}
Write \(\chi(l,s,r)=\chi_1(l)\chi_2(s)\chi_3(r)\).  If
\(\chi_3=1\), orthogonality on \(H_j\) proves the assertion when
\(\chi_2\ne1\); when \(\chi_2=1\), the average of a nontrivial
\(\chi_1\) on \(I_j\setminus\{0\}\) is eventually
\(-1/(|I_j|-1)\).

Suppose \(\chi_3\ne1\).  Fix distinct \(h_1,\ldots,h_m\in k\),
which belong to \(H_j\) for all sufficiently large \(j\).  For fixed
\(l\in I_j\setminus\{0\}\), set
\[
 f_l(s)=\chi_2(s)\chi_3((s^2+P_\theta)/l).
\]
The average of \(f_l\) equals the average of
\(m^{-1}\sum_a f_l(s+h_a)\).  In the squared modulus of the latter expression,
the correlation for \(a\ne b\) is a constant of modulus one times
the average on \(H_j\) of
\[
 s\longmapsto\chi_3\bigl(2(h_a-h_b)s/l\bigr).
\]
Choose \(z_{ab}\in k\) with
\(\chi_3(2(h_a-h_b)z_{ab})\ne1\).  Once all \(z_{ab}\) belong to
\(I_j\), the point \(lz_{ab}\in H_j\) detects this character,
uniformly for every allowed \(l\).  All these correlations are then
zero.  Cauchy--Schwarz gives
\( |\mathbb E_{s\in H_j}f_l(s)|^2\leq1/m\), uniformly in \(l\).
Letting \(m\) grow proves the character assertion.  Polynomial
avoidance follows from the product-grid assertion of
Lemma~\ref{rec:polynomial-density}.
\end{proof}

\begin{proposition}\label{rec:odd}
In odd characteristic, \(\mu\)-almost every palette satisfies
\eqref{rec:translation-invariance}.  In fact, for each
\(h\in\mathcal H\), membership in \(\mathcal S\) is constant on all
positions \(xh\), \(x\in\mathcal A\).
\end{proposition}

\begin{proof}
Fix a palette \(c\), a position \(h\), and the set
\[
 E=\{y\in\mathcal A:c_{yh}\notin\mathcal S\}.
\]
For \(x=(l_x,s_x)\notin E\), define the surjective linear map
\(L_x:k^3\to k^2\) by
\begin{equation}\label{rec:linear-map}
 L_x(l,s,r)=
 \left(\frac{P_\theta-s_x^2}{l_x}l+2s_xs-l_xr,
       2s-\frac{2s_x}{l_x}l\right).
\end{equation}
Its kernel is \(kv(x)\).  Substituting the parameters of \(yx^{-1}\)
into Equation~\ref{pal:scalar-parameters} gives exactly
\(L_xv(y)\).  First suppose \(c=c(p)\) is a rational palette, and let
\(w_x=u(xhp)\) be the unitary representative at position \(xh\).
The position \(yh\) is obtained from it by the scalar affine map
\(yx^{-1}\).  Using the map \(\mathcal B_{w_x}:k^2\to J\) of
Section~\ref{sec:palettes}, define
\[
 \nu_x=
 \bigl(\chi\longmapsto\chi\circ\mathcal B_{w_x}\circ L_x\bigr)_*
 \sigma.
\]
This positive measure on \(\widehat{k^3}\) has mass \(\delta\) and
trivial atom at least \(\delta^2\).  Its Fourier transform satisfies
\(\widehat\nu_x(v(y))=K(\mathcal B_{w_x}(L_xv(y)))\).
Proposition~\ref{pal:translation-obstruction} consequently gives
\begin{equation}\label{rec:odd-zeros}
 \widehat\nu_x(v(y))=0\quad(y\in E),\qquad
 \operatorname{supp}\nu_x\subseteq(kv(x))^\perp.
\end{equation}
The support assertion holds because every character in the image
annihilates \(\ker L_x\).

These measures exist for arbitrary palettes as well.  Approximate \(c\)
by rational palettes on successively larger finite sets of positions.
For each fixed \(x\notin E\), take a weak limit of the corresponding
measures on the compact dual.  The equations in
\eqref{rec:odd-zeros} pass to the limit, as does support in the closed
annihilator \((kv(x))^\perp\).  The lower bound on the trivial atom
persists because \(\{1\}\) is closed.

We claim that either \(E=\mathcal A\), or \(E\) has density zero.
By Lemmas~\ref{rec:spectral-concentration} and
\ref{rec:quadric-orthogonality}, a nonempty complement has lower density
at least \(\delta\).  If \(E\) had positive density along a
subsequence, those lemmas would give a finite set of nontrivial
characters such that, for every \(x\notin E\), at least one of them
annihilates \(kv(x)\).

For any fixed nontrivial \(\chi\), this annihilation condition is a
proper polynomial condition on \((l_x,s_x)\).  To treat arbitrary
additive characters, fix a primitive
\(p\)-th root of unity and regard characters as
\(\bbF_p\)-linear functionals.  The space
\(\Hom_{\bbF_p}(k,\bbF_p)\) is a \(k\)-vector space under
\[
 (a\lambda)(t)=\lambda(at).
\]
If \(\chi\) has component functionals \(\lambda_1,\lambda_2,\lambda_3\),
annihilation of \(kv(l,s)\) says
\[
 l\lambda_1+s\lambda_2+\frac{s^2+P_\theta}{l}\lambda_3=0.
\]
This is a rational vector equation in the finite-dimensional span of
the three functionals.  It is not an identity: the vectors \(v(l,s)\)
span \(k^3\), as follows on multiplying any putative linear relation
among their coordinates by \(l\).  Clearing denominators therefore
places its solutions in a proper polynomial locus.  A finite union of
these loci has density zero, contradicting the positive lower density
of \(\mathcal A\setminus E\).  The claim follows.

It remains to turn this density dichotomy into actual invariance.
Let \(q_0=|\mathcal F\setminus\mathcal S|/|\mathcal F|\).  For every
fixed \(h\), uniform marginals give
\(\int\mathbb E_j\boldsymbol1_E\,d\mu=q_0\).  The claim and dominated
convergence show that \(\mu(E=\mathcal A)=q_0\).  For each
\(x\in\mathcal A\), the event \(E=\mathcal A\) is contained in
\(\{c_{xh}\notin\mathcal S\}\), whose probability is also \(q_0\).
Their difference is null.  There are only countably many \(x\) and
\(h\), so outside a single null set membership is constant on every
one of the asserted families of positions.
\end{proof}

\subsection{Characteristic two}

Assume now that \(\operatorname{char}k=2\), so \(T\ne0\).  We shall
prove scalar-translation invariance directly on rational points.  A
unitary element \(w\) belongs to the field \(E=L[w]\subset D\), which
is stable under \(*\) because \(w^*=w^{-1}\).  With
\(F=E^{\langle *\rangle}\), we
have \(E=LF\).  The degree \([E:L]\) divides the odd degree of \(D\),
so \(E/L\), and hence \(F/k\), is separable.  The Cayley chart identifies
\((E/F)^1\) with \(F\cup\{\infty\}\).

We first prove a density statement on each of these fields.  Throughout
this subsection, densities on \(k\) are taken over an arbitrary fixed
increasing sequence of finite additive subgroups exhausting \(k\).

\begin{lemma}\label{rec:two-density}
Let \(F\subset D\) be as above, let \(Y\in F\), and let
\(H_0\in k[X]\) be a nonconstant additive polynomial.  For
\[
 w_x=q(Y+H_0(x)),\qquad x\in k,
\]
either \(\phi(w_x)\notin\mathcal S\) for every \(x\), or the set of
such \(x\) has density zero.
\end{lemma}

\begin{proof}
Put \(E_0=\{x:\phi(w_x)\notin\mathcal S\}\) and
\(Q_0=H_0^2+TH_0\).  The polynomial \(Q_0\) is nonconstant and
additive.  Equation~\ref{pal:scalar-parameters}, applied to scalar
translations, gives
\begin{equation}\label{rec:two-zero}
 K\bigl(B(x)Q_0(y-x)\bigr)=0\qquad(x\notin E_0,\ y\in E_0),
\end{equation}
where
\begin{equation}\label{rec:two-coefficient}
 B(x)=\frac{w_x+w_x^*}{T}
     =\frac{T}{(Y+H_0(x))^2+T(Y+H_0(x))+P_\theta}\in J.
\end{equation}
In any fixed \(k\)-basis of \(J\), the coordinates of \(B(x)\) are
rational functions of \(x\) tending to zero at infinity.  This follows
also directly by expanding the inverse denominator in
\(F((x^{-1}))\).  The denominator never vanishes for \(x\in k\),
since its two factors are \(Y+H_0(x)+\theta\) and
\(Y+H_0(x)+\bar\theta\), and \(\theta\notin F\).

Push \(\sigma\) forward by the character map
\[
 \chi\longmapsto
 \bigl[y\longmapsto\chi(B(x)Q_0(y))\bigr]
\]
to obtain a measure \(\nu_x\) on \(\widehat k\).  It has mass
\(\delta\) and trivial atom at least \(\delta^2\).
Equation~\eqref{rec:two-zero} is precisely the vanishing required by
Lemma~\ref{rec:spectral-concentration}, with \(v(y)=y\) and \(a_x=x\).
Ordinary orthogonality on exhausting additive subgroups supplies
\eqref{rec:orthogonality}.  Thus, if \(k\setminus E_0\ne\varnothing\),
it has positive lower density.  If \(E_0\) also had density tending to
some \(d>0\) along a subsequence, there would be a fixed finite set
\(\Lambda\subset\widehat k\setminus\{1\}\) such that
\begin{equation}\label{rec:two-spectral-mass}
 \sigma\{\chi:\chi(B(x)Q_0(\cdot))\in\Lambda\}
       \geq d\delta^2/2\qquad(x\notin E_0).
\end{equation}

We prove that, for every fixed \(\chi\in\widehat J\) and every fixed
nontrivial \(\lambda\in\widehat k\), the equality
\begin{equation}\label{rec:wild-character}
 \chi(B(x)Q_0(\cdot))=\lambda
\end{equation}
holds on a set of density zero.  This is the point where the additive
polynomial and the use of arbitrary additive characters both matter.
Regard the characters as \(\bbF_2\)-linear functionals, with the
\(k\)-vector space structure used in the preceding subsection.  Write
\(Q_0(y)=\sum_i a_i y^{2^i}\), and choose \(M=2^r\geq\deg Q_0\).
The extension \(k/k^M\) is finite.  For a basis \(e_1,\ldots,e_N\)
over \(k^M\), the representation
\[
 x=\sum_{j=1}^N e_jx_j^M
\]
is unique and defines an additive group isomorphism \(k^N\to k\).
Let \(b_\ell(x)\) be the coordinates of \(B(x)\), and
\(\chi_\ell\) the component functionals of \(\chi\).  For each
\(i\), choose a basis \((f_{im})_m\) of \(k/k^{2^i}\).  Expansion in
that basis gives rational functions \(R_{\ell im}\in k(X_1,\ldots,X_N)\)
such that
\begin{equation}\label{rec:frobenius-expansion}
 a_i b_\ell\left(\sum_j e_jX_j^M\right)
       =\sum_m f_{im}R_{\ell im}(X)^{2^i}.
\end{equation}
Indeed the left side belongs to
\(k(X_1^{2^i},\ldots,X_N^{2^i})\), and this field has basis
\((f_{im})_m\) over \(k^{2^i}(X_1^{2^i},\ldots,X_N^{2^i})\).

Define fixed functionals
\(\psi_{\ell im}(y)=\chi_\ell(f_{im}y^{2^i})\).  The left side of
\eqref{rec:wild-character}, in functional notation, is now
\begin{equation}\label{rec:functional-rational-map}
 \sum_{\ell,i,m}R_{\ell im}(X)\psi_{\ell im}.
\end{equation}
It is a rational map into a finite-dimensional \(k\)-vector space.
Moreover, it tends to zero when all \(X_j\) are multiplied by a common
indeterminate tending to infinity.  For completeness, the left side of
\eqref{rec:frobenius-expansion} tends to zero under this scaling because
each \(b_\ell\) vanishes at infinity.  If some
\(R_{\ell im}\) did not tend to zero, the terms of largest Laurent
degree on the right could not cancel: the \(f_{im}\) are linearly
independent over
\(k^{2^i}(X_1^{2^i},\ldots,X_N^{2^i})\).  Thus every one of the
rational functions in \eqref{rec:functional-rational-map} tends to
zero.

It follows that this rational map cannot be identically equal to the
fixed nonzero functional \(\lambda\).  If \(\lambda\) is outside its
finite-dimensional span the equality is impossible; otherwise, taking
coordinates and clearing denominators puts its solutions in a proper
polynomial locus.  Any exceptional denominator locus is proper as well.
Lemma~\ref{rec:polynomial-density}, transported through the additive
isomorphism \(k^N\to k\), proves the asserted density zero.

Average \eqref{rec:two-spectral-mass} over \(x\).  Its right side,
restricted to \(k\setminus E_0\), has positive lower average.  Its left
side has average tending to zero: for each fixed \(\chi\), the finite
union of the sets \eqref{rec:wild-character}, with
\(\lambda\in\Lambda\), has density zero, and dominated convergence
applies to the finite measure \(\sigma\).  This contradiction proves
the density dichotomy.
\end{proof}

The density statement alone allows exceptional points.  The abelian
structure of the colors on a field removes them.

\begin{proposition}\label{rec:two}
In characteristic two, membership in \(\mathcal S\) is unchanged by
every scalar translation on \(X(k)\).  Consequently every palette
satisfies \eqref{rec:translation-invariance}.
\end{proposition}

\begin{proof}
Fix \(E=LF\) as above and let
\(A_E=\phi((E/F)^1)\), a finite abelian subgroup of \(\mathcal F\).
If \(A_E\cap\mathcal S\) is empty there is nothing to prove.  Otherwise
Theorem~\ref{fin:obstruction} says that
\(A_E\setminus\mathcal S\) is a subgroup, say \(A_E^0\).  Define the
homomorphism, with additive target notation,
\[
 \psi:E^\times\longrightarrow A_E/A_E^0,
 \qquad \psi(a)=\phi(a/\bar a)\bmod A_E^0,
\]
where bar denotes the involution on \(E\).  Membership of
\(\phi(q(Z))\) in \(\mathcal S\) is equivalent to
\(\psi(Z+\theta)\ne0\).

For \(Y\in F\), consider the finite-valued function
\(g(s)=\psi(Y+s+\theta)\) on \(k\).  We show that its additive period
group
\[
 \Pi=\{t\in k:g(s+t)=g(s)\text{ for every }s\in k\}
\]
has finite index.  For \(t\ne0\), agreement of these two values is
equivalent to vanishing of \(\psi\) on
\begin{align*}
 &(Y+s+t+\theta)(Y+s+\bar\theta)\\
 &\hspace{1cm}=Y^2+(T+t)Y+P_\theta+s^2+(T+t)s+t\bar\theta.
\end{align*}
Here \(\psi(\bar a)=-\psi(a)\), so the product measures the difference
of the two values of \(g\).  Division by \(t\in k^\times\) does not
change \(\psi\).  With
\[
 Z_t=\frac{Y^2+(T+t)Y+P_\theta}{t},\qquad
 H_t(s)=\frac{s^2+(T+t)s}{t},
\]
the agreement condition becomes
\(\psi(Z_t+H_t(s)+\bar\theta)=0\).  Replacing \(\bar\theta\) by
\(\theta\) negates this value and preserves its zero set.  The
polynomial \(H_t\) is nonconstant and additive, so
Lemma~\ref{rec:two-density} says that either agreement holds for every
\(s\), or its set has density zero.  Thus every \(t\notin\Pi\)
gives density-zero agreement.

Choose translations of \(g\) indexed by distinct cosets of \(\Pi\).
Any two agree on a set of density zero, because the averages are
eventually invariant under each fixed additive shift.  If there were
more such translates than values of \(g\), then at every point two
would agree.  A finite union of density-zero sets cannot cover \(k\).
This proves that \([k:\Pi]\) is finite.

Finally apply Lemma~\ref{rec:two-density} with \(H_0(s)=s\).  The zero
set of \(g\) is either all of \(k\) or has density zero.  Since \(g\)
is constant on cosets of \(\Pi\), any nonempty zero set contains a
coset of density \([k:\Pi]^{-1}\).  Hence \(g\) is either identically
zero or nowhere zero.  Membership is therefore constant on
\(\{q(Y+s):s\in k\}\).  Every finite rational point belongs to such
a field chart, and \(\infty\) is fixed by translations.  This proves
the rational-point assertion.  For fixed \(s\in k\) and
\(h\in\mathcal H\), the equality of membership indicators
\(\boldsymbol1_{\mathcal S}(c_{\tau_s h})=
\boldsymbol1_{\mathcal S}(c_h)\) is a closed condition.  It holds for
every rational palette, hence for every palette in their closure.
\end{proof}

\subsection{Generating rotations}

In either characteristic, choose a palette \(c\) satisfying
\eqref{rec:translation-invariance}.  Define
\[
 \mathcal G_c=\{g\in\mathcal H:I_c(gh)=I_c(h)
                       \text{ for every }h\in\mathcal H\}.
\]
This is a subgroup.  In addition to all \(\tau_s\), it contains left
rotations \(L_a\) with \(\phi(a)=1\), and all unitary conjugations:
these assertions follow from the palette rotation identities and the
conjugacy invariance of \(\mathcal S\).  We now explain how these
transformations generate squares of arbitrary separable field rotations.
The construction adapts \cite[Section~5.6]{MPNF}, with the field-spanning
step made explicit in positive characteristic.

\begin{lemma}\label{rec:field-rotations}
Let \(E=LF\subset D\) be a field stable under \(*\), with
\(F=E^{\langle *\rangle}\) and \(E/L\) separable.  Then
\(L_{\xi^2}\in\mathcal G_c\) for every \(\xi\in(E/F)^1\).
\end{lemma}

\begin{proof}
We first obtain all \(F\)-translations, then use a central rotation to
generate \(\SL_2(F)\), and finally express a square rotation in Cayley
coordinates.  Choose
\(n_0=q(x_0)\in(E/F)^1\cap\ker\phi\), with \(x_0\) generating
\(F/k\), such that the polynomials
\begin{equation}\label{rec:disjoint-poles}
 (s+x_i+\theta)(s+x_i+\bar\theta)
\end{equation}
have pairwise disjoint root sets as \(x_i\) runs over the conjugates
of \(x_0\) under embeddings fixing \(L\).  These are nonempty Zariski
open conditions: over a splitting field the conjugate coordinates of
\(\Res_{F/k}\bbA^1\) are independent, and all prohibited equalities
are proper linear conditions.  They can be met inside \(\ker\phi\)
because powers annihilating \(\mathcal F\) are Zariski dense in the
norm-one torus.  Indeed its rational points are Zariski dense by the
Cayley parametrization, and every positive power map is dominant, even
when it is inseparable.

Since the color set is finite, there are \(t_0\in k\) and an infinite
set \(\Sigma\subset k\) for which
\(\phi(\tau_s(n_0))=\phi(\tau_{t_0}(n_0))\) for every
\(s\in\Sigma\).  Here translations on unitary elements are understood
through the Cayley chart.  Put
\[
 g_s=L_{n_0}^{-1}\tau_s^{-1}
       L_{\tau_s(n_0)\tau_{t_0}(n_0)^{-1}}
       \tau_{t_0}L_{n_0}\in\mathcal G_c.
\]
Each \(g_s\) fixes the unitary representative \(1\), which is the
point \(\infty\) in additive coordinates.  Its fractional matrix has
entries in \(E\) and is upper triangular.  To compute its diagonal
entries, the exact matrix of \(L_{q(z)}\), for \(z\in F\), is
\[
 M_z=(z+\bar\theta)^{-1}
       \begin{pmatrix}z&-P_\theta\\1&z+T\end{pmatrix}.
\]
The factor \((z+\bar\theta)^{-1}\) is retained here because it need
not be central in \(D\).  Multiplying the matrices in the expression
for \(g_s\) gives diagonal entries
\[
 a_s=\frac{x_0+t_0+\bar\theta}{x_0+s+\bar\theta},\qquad
 d_s=\frac{x_0+s+\theta}{x_0+t_0+\theta}.
\]
Thus its affine multiplier on the field line
\(F\cup\{\infty\}\) is
\begin{equation}\label{rec:affine-multipliers}
 m_s=\frac{(x_0+t_0+\theta)(x_0+t_0+\bar\theta)}
           {(x_0+s+\theta)(x_0+s+\bar\theta)}\in F^\times.
\end{equation}
This is also the ratio of the logarithmic derivatives of \(q\) at
\(x_0+s\) and \(x_0+t_0\).

The same multiplier controls conjugation of scalar translations on the
whole chart \(J\).  Indeed an upper triangular fractional matrix
\(\left(\begin{smallmatrix}a&b\\0&d\end{smallmatrix}\right)\)
acts by \(x\mapsto(ax+b)d^{-1}\).  Conjugating \(x\mapsto x+r\),
\(r\in k\), gives \(x\mapsto x+r ad^{-1}\).  Here
\(ad^{-1}=m_s\), so \(\mathcal G_c\) contains translations by
\(km_s\).

The elements \(m_s\), \(s\in\Sigma\), span \(F\) over \(k\).
Otherwise a nonzero \(k\)-linear functional on \(F\) would annihilate
them.  After passing to a normal closure, such a functional is a linear
combination of the distinct embeddings of \(F\).  It would give a
linear relation among the rational functions
\[
 s\longmapsto
 \frac{(x_i+t_0+\theta)(x_i+t_0+\bar\theta)}
      {(x_i+s+\theta)(x_i+s+\bar\theta)}
\]
at infinitely many \(s\), and therefore an identity.  Each summand
has a pole that occurs in no other summand by
\eqref{rec:disjoint-poles}; its numerator is nonzero.  Every coefficient
must consequently vanish, a contradiction.  Composing the translations
already obtained now gives every \(F\)-translation on \(J\).

For a central rotation, take the particular scalar
\(a=q(0)=\theta/\bar\theta\ne1\).  Its color is trivial, so
\(L_a\in\mathcal G_c\).  In additive coordinates it acts by
\(x\mapsto-P_\theta(x+T)^{-1}\).  Consequently
\(W=L_a\tau_{-T}\) is represented by
\(\left(\begin{smallmatrix}0&-P_\theta\\1&0\end{smallmatrix}\right)\),
and direct multiplication gives
\[
 W\begin{pmatrix}1&b\\0&1\end{pmatrix}W^{-1}
       =\begin{pmatrix}1&0\\-b/P_\theta&1\end{pmatrix}
       \qquad(b\in F).
\]
Thus all lower as well as upper unipotent matrices over \(F\) act in
\(\mathcal G_c\).  Hence \(\mathcal G_c\) contains the fractional action
on \(J\) of \(\SL_2(F)\), since upper and lower elementary matrices
generate that group.

Finally set
\[
 C=\begin{pmatrix}1&\theta\\1&\bar\theta\end{pmatrix}.
\]
For \(\xi\in(E/F)^1\), the matrix
\begin{equation}\label{rec:cayley-diagonal}
 C^{-1}\begin{pmatrix}\xi&0\\0&\xi^{-1}\end{pmatrix}C
 =\frac1{\bar\theta-\theta}
 \begin{pmatrix}
 \bar\theta\xi-\theta\xi^{-1}&P_\theta(\xi-\xi^{-1})\\
 \xi^{-1}-\xi&\bar\theta\xi^{-1}-\theta\xi
 \end{pmatrix}
\end{equation}
has all entries in \(F\) and determinant one.  It belongs to the
\(\SL_2(F)\) just generated.  In unitary coordinates its action is
\(u\mapsto\xi u\xi\), as is seen by applying the diagonal matrix to
the column \((x+\theta,x+\bar\theta)^{\mathsf t}\).  Composing this
action with the unitary conjugation \(u\mapsto\xi u\xi^{-1}\)
gives \(u\mapsto\xi^2u\).  This is the required left rotation.
\end{proof}

\begin{proposition}\label{rec:no-simple-colors}
There is no homomorphism \(\phi:U(k)\to\mathcal F\) to a finite
nonabelian simple group that kills \(R\) and satisfies
\(\phi(V_0)=\mathcal F\).
\end{proposition}

\begin{proof}
Every color is represented by an element \(\xi\in V_0\) lying in a
separable field of the kind used in Lemma~\ref{rec:field-rotations}.
Indeed, start with any representative and perturb it within its
\(R\)-coset to a regular semisimple element, using
Proposition~\ref{arith:closure} and the nonempty local regular semisimple
open subset.  In a division algebra its generated algebra over \(L\)
is a field, stable under the involution, and regular semisimplicity
makes that field separable.  This also covers characteristics dividing
the odd degree \(n\).

Lemma~\ref{rec:field-rotations} and the left-rotation palette rule
therefore make the membership indicator invariant under left
multiplication by the square of every element of \(\mathcal F\).
These squares generate \(\mathcal F\): their generated subgroup is
normal, and if it were trivial then \(\mathcal F\) would have exponent
two and hence be abelian.  Since \(\phi\) is surjective, the positions
obtained from any one position by left rotations realize every color.
The indicator must be constant on \(\mathcal F\), contrary to the
choice of the nonempty proper set \(\mathcal S\).
\end{proof}
